PST / OPP – Numerische und theoretische Untersuchungen

Gesamtwerk – Inhaltsverzeichnis

Band I – Grundlagen und Motivation

1. Einleitung
2. Motivation der Projective Structure Theory (PST)
3. Der Ontophänomenale Phasenraum (OPP)
4. Grundannahmen des Ontophänomenalismus
5. Ziel der numerischen Untersuchungen
6. Offene Forschungsfragen

---

Band II – Mathematische Grundlagen

1. Relationale Kernel
2. Konstruktion der Gram-Matrix
3. Mutual Information
4. Spektralzerlegung
5. Effektive Dimension
6. Partizipationszahl
7. Stabilitätsfunktional
8. Optimierungsverfahren
9. Projektionen
10. Diffusionsoperatoren
11. Kernel-PCA
12. Random Projection
13. Diffusion Maps
14. Mathematische Interpretation

---

Band III – Numerische Simulationen

Kapitel 1 – Erste Untersuchungen

- Test 022
  Selbstkonsistente Projektion

- Test 023
  Abhängigkeit vom Zielrang

- Test 024
  weitere Kontrollsimulationen

- Test 025
  Optimierung des universellen Stabilitätsfunktionals

- Test 026
  Robustheitstest

- Test 027
  Zielrang-Analyse

- Test 028
  Spektralform-Audit

- Test 029
  Spectral-Gap-Analyse

---

Kapitel 2 – Projektionsmethoden

- Test 030
  Kernel-Universalität

- Test 031
  Rohskalierung des OPP

- Test 031B
  Skalierung der effektiven Dimension

- Test 032
  PCA gegen Diffusion Maps

- Test 033
  Kernel PCA gegen Random Projection

---

Kapitel 3 – Diffusionsuntersuchungen

- Test 034
  Diffusionszeit-Scan

- Test 035
  Feinscan der Diffusionszeit

- Test 036
  Größenabhängigkeit

- Test 037
  Universelles Skalierungsgesetz

---

Band IV – Zusammenfassung aller numerischen Ergebnisse

1. Was ist sicher nachgewiesen?
2. Welche Resultate sind robust?
3. Welche Resultate hängen von Parametern ab?
4. Welche Resultate sprechen gegen die PST?
5. Welche Resultate sprechen für die PST?
6. Welche Fragen bleiben offen?

---

Band V – Physikalische Interpretation

1. Emergenz der Raumzeit
2. Beobachter und Projektion
3. Informationsgeometrie
4. Diffusionsgeometrie
5. Kernel-Geometrie
6. Projektionsprinzip
7. Lorentz-Invarianz
8. Quantisierung
9. Begrenzte Auflösung des Beobachters
10. Zusammenhang mit Eichsymmetrien
11. Lie-Gruppen
12. Kaluza-Klein
13. Höherdimensionale Strukturen
14. Projektive Raumzeit

---

Band VI – Ontophänomenalismus

1. Existenz
2. Bestimmtheit
3. Differenz
4. Relation
5. Objekt
6. Relative Existenz (RdE)
7. Projektion
8. Spurargument
9. Vergangenheit und Zukunft
10. Beobachter
11. Verbindung zwischen OP und PST

---

Band VII – Offene Probleme

1. Warum genau diese Projektion?
2. Existiert eine universelle Projektion?
3. Entsteht Lorentz-Geometrie automatisch?
4. Entstehen Eichgruppen automatisch?
5. Ist Raumzeit vollständig emergent?
6. Ist die Quantisierung rein projektiv?
7. Entsteht Gravitation zwangsläufig?
8. Wie lässt sich das experimentell testen?

---

Band VIII – Ausblick

1. Neue numerische Experimente
2. Erweiterungen des OPP
3. Alternative Kernel
4. Komplexe Relationen
5. Verbindung zur Quanteninformation
6. Verbindung zur Allgemeinen Relativität
7. Verbindung zum Standardmodell
8. Langfristige Forschungsziele

---

Anhang

A. Vollständige Formelsammlung

B. Definition aller verwendeten Größen

C. Tabellen sämtlicher numerischer Ergebnisse

D. Glossar

E. Literatur

F. Chronologie der Entwicklung der PST

G. Chronologie aller numerischen Tests

H. Offene mathematische Vermutungen
